\section*{Roadmap}
\addcontentsline{toc}{section}{Roadmap}

\noindent\motiv{\revs{The chapters are ordered by logical dependence, and Figure~\ref{fig:roadmap}
shows which chapters use which.}}

\begin{figure}[ht]
\centering
\begin{tikzpicture}[
  font=\small,
  done/.style={draw, rounded corners=2pt, fill=Blue!6, text width=4cm, align=center, execute at begin node={\hyphenpenalty=10000\relax},
    minimum height=1.05cm, inner sep=3pt},
  plan/.style={done, dashed, fill=none},
  dep/.style={-{Stealth[length=2mm]}, thick},
  node distance=0.9cm and 0.7cm]
  \node[done] (c1) {\hyperref[sec:gaussianqca]{Q1\quad Translation-invariant maps and their symbols}};
  \node[done, right=of c1] (c2) {\hyperref[sec:ccrautomata]{Q2\quad Gaussian cellular automata}};
  \node[done, right=of c2] (c3) {\hyperref[sec:gaussianstates]{Q3\quad Gaussian states on the lattice}};
  \node[done, below=of c3] (c4)
    {\hyperref[sec:stationary]{Q4a\quad The example system and its stationary states}};
  \node[done, below=of c4] (c4b) {\hyperref[sec:convergence]{Q4b\quad Convergence to the stationary states}};
  \node[done, below=of c2] (c5) {\hyperref[sec:channels]{Q5a\quad Gaussian channels}};
  \node[done, below=of c5] (c5b)
    {\hyperref[sec:chain]{Q5b\quad Irreversible Gaussian automata on the chain}};
  \draw[dep] (c1) -- (c2);
  \draw[dep] (c2) -- (c3);
  \draw[dep] (c3) -- (c4);
  \draw[dep] (c4) -- (c4b);
  \draw[dep] (c2) -- (c5);
  \draw[dep] (c3) -- (c5);
  \draw[dep] (c5) -- (c5b);
  \draw[dep] (c4) -- (c5b);
\end{tikzpicture}
\caption{Milestones and their dependence. Each box links to its chapter.}
\label{fig:roadmap}
\end{figure}
